%% toc: P3 40 40.9
\chapter{Ehrenfeucht-Game Proofs of Asymptotic Results}
\label{sec:p03-w03-p03-072-game-proofs}

A zero--one law for a random structure says that every first-order sentence
has limiting probability \(0\) or \(1\).  Spencer's route to every such law
has two halves: a probabilistic half, which shows that each axiom of an
explicit theory \(T\) holds with probability tending to one, and a logical
half, which shows that \(T\) is complete.  This chapter is about the logical
half when it is carried out with the Ehrenfeucht game
\(\mathrm{EHR}(G,H;t)\), and about the step that joins the two halves.  It
proves the game theorem in the form Spencer states it, marks the two
hypotheses of his generalization and of his completeness criterion that
cannot be dropped, proves that almost sure axioms and completeness give a
zero--one law, and assembles Duplicator's closure strategy for the sparse
random graph \(G(n,n^{-\alpha})\) into a winning strategy for the whole game.

\emph{Inputs.} A first-order language with equality and finitely many
relation symbols; two structures for it; a number of rounds \(t\); for the
asymptotic applications, a probability space over the structures with
domain \(\{1,\ldots,n\}\) for each large \(n\) and a candidate theory \(T\).
For the sparse random graph, an exponent \(\alpha\in(0,1)\) and the
weighted notions of rigidity, safety and closure of
\cref{def:p03-w03-p03-070-weight,def:p03-w03-p03-070-closure}.

\emph{Conclusions.} Duplicator wins \(\mathrm{EHR}(G,H;t)\) exactly when
\(G\) and \(H\) agree on the sentences of quantifier depth \(t\), in either
reading of that phrase (\cref{thm:p03-w03-p03-072-game-theorem}); with
infinitely many relation symbols this fails
(\cref{thm:p03-w03-p03-072-infinite-language}); a complete theory of almost
sure sentences gives a zero--one law whose limit-one sentences are its
consequences (\cref{thm:p03-w03-p03-072-transfer}); and on any two models
of the almost sure theory \(T_\alpha\), \(\alpha\) irrational, Duplicator
wins every finite game by keeping closures isomorphic
(\cref{thm:p03-w03-p03-072-closure-strategy}).

\emph{What is treated elsewhere.} Spencer's account of the game itself,
the remark that Spoiler may be required to choose new vertices, his
isolated-vertex example, the elementary-equivalence and completeness
corollaries and the gap between all finite games and the infinite game are
in \cref{sec:p03-w03-p03-064-game}
(\cref{thm:p03-w03-p03-064-new-vertices,def:p03-w03-p03-064-isolated-game,thm:p03-w03-p03-064-ee-game,thm:p03-w03-p03-064-complete-game,thm:p03-w03-p03-064-infinite-game}).
The closure, the Nonexistence and Finite Closure Lemmas and the axioms of
\(T_\alpha\) are in \cref{def:p03-w03-p03-064-closure,thm:p03-w03-p03-068-finite-closure,def:p03-w03-p03-070-theory};
the single round of the closure strategy is proved, for a pair of
parameters, as \cref{thm:p03-w03-p03-070-extension}.  This chapter uses
those results and does not restate them.

\emph{Version followed.} The presentation follows the separately formatted
version of ``Logic and Random Structures'' that Spencer posted on his web
page \cite{Spencer2007LogicRandomStructuresAuthor}, a
30-page text with sections numbered 1 to 14 and its own page numbers, which
are the page numbers cited here.  The game, its theorem, the two
corollaries and the extension to other languages are in its Section~14
\cite[pp.~26--27]{Spencer2007LogicRandomStructuresAuthor}; the closure
strategy is in its Section~5
\cite[pp.~11--12]{Spencer2007LogicRandomStructuresAuthor}.  The chapter as
printed in the Springer volume \emph{Finite Model Theory and Its
Applications} (2007) was not consulted, and no locator here refers to it.

\section{The game theorem in Spencer's form}
\label{sec:p03-w03-p03-072-theorem}

Spencer defines \(\mathrm{EHR}(G,H;t)\) first for graphs and then, with the
same board, players and play, for any language with finitely many relation
symbols of various arities: Duplicator must then preserve every relation,
that is, \(R(x_{i_1},\ldots,x_{i_l})\) and \(R(y_{i_1},\ldots,y_{i_l})\)
must have the same truth value for every \(l\)-ary \(R\) and all indices
\(i_1,\ldots,i_l\) from \(1,\ldots,t\), besides \(x_i=x_j\) iff
\(y_i=y_j\) \cite[pp.~26--27]{Spencer2007LogicRandomStructuresAuthor}.
For graphs this is \cref{def:p03-w03-p03-067-game}; in general it is the
game of \cref{def:p03-ef-game} for a finite relational signature.  Spencer
writes the game both as \(\mathrm{EHR}(G,H;t)\) and as
\(\mathrm{EHR}[G,H;t]\).  He calls the next statement fundamental and gives
it without proof \cite[pp.~26--27]{Spencer2007LogicRandomStructuresAuthor}.
The Ehrenfeucht--Fra\"iss\'e theorem is recorded in
\cref{thm:p03-ef} with a citation; the handbook proves it, with parameters,
as \cref{thm:p03-w03-p03-038-pointed-ef}.  The annex supplies the bridge
from Spencer's game, and from his phrase ``quantifier depth \(t\)'', to that
proved theorem.  Quantifier depth is the quantifier rank
\(\operatorname{qr}\) of \cref{def:p03-quantifier-rank}.

\begin{hbtheorem}{thm:p03-w03-p03-072-game-theorem}{proved=annex:p03-w03-p03-072-game-theorem}{The game theorem}
Let \(L\) be a first-order language with equality, finitely many relation
symbols and no function or constant symbols.  Let \(G,H\) be
\(L\)-structures with nonempty domains, and let \(t\geq1\).  The following
are equivalent.
\begin{enumerate}
\item Duplicator wins \(\mathrm{EHR}(G,H;t)\).
\item \(G\) and \(H\) give the same truth value to every \(L\)-sentence of
  quantifier depth exactly \(t\).
\item \(G\) and \(H\) give the same truth value to every \(L\)-sentence of
  quantifier depth at most \(t\).
\end{enumerate}
\end{hbtheorem}

The equivalence of (2) and (3) uses that the domains are nonempty: a
sentence of smaller depth is padded to depth \(t\) by a conjunct that is
true in every nonempty structure.  Spencer's illustration, a graph with an
isolated vertex against one without, is
\cref{def:p03-w03-p03-064-isolated-game}.

Spencer's generalization is to languages with a \emph{finite} number of
relation symbols \cite[p.~27]{Spencer2007LogicRandomStructuresAuthor}.  The
restriction is needed for the direction from (3) to (1), as the next
example, which is the handbook's and not the source's, shows.

\begin{hbtheorem}{thm:p03-w03-p03-072-infinite-language}{proved=annex:p03-w03-p03-072-infinite-language}{With infinitely many relation symbols the game theorem fails}
Let \(L\) have equality and unary relation symbols \(U_1,U_2,\ldots\).  Let
\(G\) have domain \(\{a_1,a_2,\ldots\}\) with \(G\entails U_k(a_n)\) iff
\(k\leq n\), and let \(H\) extend \(G\) by one further element
\(a_\infty\) with \(H\entails U_k(a_\infty)\) for every \(k\).  Then \(G\)
and \(H\) agree on every \(L\)-sentence of quantifier depth at most \(1\),
and Spoiler wins \(\mathrm{EHR}(G,H;1)\) played with Spencer's rules.
\end{hbtheorem}

The proof of the theorem for finite languages rests on there being only
finitely many formulas of each depth in a fixed finite set of variables, up
to equivalence.  In the example the single element \(a_\infty\) has a type
that no element of \(G\) has, although each finite part of that type is
realized in \(G\).

\begin{hbwarning}{def:p03-w03-p03-072-finite-models}{The completeness criterion needs a theory without finite models}
Spencer's criterion, \cref{thm:p03-w03-p03-064-complete-game}, assumes that
the consistent theory \(T\) has no finite models
\cite[p.~27]{Spencer2007LogicRandomStructuresAuthor}.  Its ``if'' direction
fails without that hypothesis when ``countable model'' is read as
``countably infinite model''.  Let \(T\) say that a graph has no edges and
at most two vertices:
\(\forall x\,\forall y\;\neg x\sim y\) and
\(\forall x\,\forall y\,\forall z\,(x=y\lor x=z\lor y=z)\).  \(T\) is
consistent and has no infinite model, so the game condition holds
vacuously.  Yet \(T\) is not complete: \(\exists x\,\exists y\;x\neq y\)
holds in its two-vertex model and fails in its one-vertex model.  The
hypothesis is what makes the countable models of the criterion infinite,
and its proof obtains them from arbitrary models by the downward
L\"owenheim--Skolem theorem.  In the asymptotic applications the hypothesis
holds automatically: an almost sure theory contains, for every \(r\), the
sentence saying that there are \(r\) distinct elements, since it holds in
every structure of size at least \(r\)
\cite[p.~13]{Spencer2007LogicRandomStructuresAuthor}
(\cref{thm:p03-w03-p03-067-as-theory}).
\end{hbwarning}

\section{From almost sure axioms to a zero--one law}
\label{sec:p03-w03-p03-072-transfer}

Spencer's opening example sets out the pattern that all of his zero--one
laws follow: find an explicit theory \(T\) all of whose members hold almost
surely, and prove \(T\) complete
\cite[p.~3]{Spencer2007LogicRandomStructuresAuthor}.  He observes that the
notion of almost sure truth makes sense whenever there is, for every large
\(n\), a probability space over models of size \(n\)
\cite[p.~2]{Spencer2007LogicRandomStructuresAuthor}; the step joining the
two halves holds at that generality.

\begin{hbconvention}{def:p03-w03-p03-072-random-structure}{Random structures and almost sure sentences}
Fix a first-order language \(L\) with equality, finitely many relation
symbols and no function or constant symbols.  A \emph{random
\(L\)-structure} assigns to every sufficiently large \(n\) a probability
measure \(\Pr_n\) on the finite set of \(L\)-structures with domain
\(\{1,\ldots,n\}\).  An \(L\)-sentence \(A\) holds \emph{almost surely} if
\(\lim_{n\to\infty}\Pr_n[A]=1\) and \emph{almost never} if this limit is
\(0\), equivalently if \(\neg A\) holds almost surely
\cite[p.~2]{Spencer2007LogicRandomStructuresAuthor}.  For graphs these are
the notions of \cref{def:p03-w03-p03-065-almost-sure}; the unary model
\(SU(n,p)\) and the random graph \(G(n,p)\), with \(p=p(n)\) a function of
\(n\), are random structures in this sense
(\cref{def:p03-w03-p03-067-models,def:p03-w03-p03-065-gnp}).
\end{hbconvention}

\begin{hbtheorem}{thm:p03-w03-p03-072-transfer}{proved=annex:p03-w03-p03-072-transfer}{Almost sure and complete gives zero--one}
Let a random \(L\)-structure be given, and let \(T\) be a set of
\(L\)-sentences each of which holds almost surely.
\begin{enumerate}
\item Every consequence of \(T\) holds almost surely, and \(T\) is
  consistent.
\item If \(T\) is complete, then every \(L\)-sentence \(B\) holds almost
  surely or almost never; \(B\) holds almost surely iff \(T\entails B\); and
  the set of almost sure sentences is the deductive closure of \(T\).
\end{enumerate}
\end{hbtheorem}

Item 1 is Spencer's ``finiteness of proof'' argument and item 2 his
appeal to completeness \cite[p.~3]{Spencer2007LogicRandomStructuresAuthor}.
The converse, that the almost sure sentences of a zero--one law form a
complete theory with countably infinite models, is
\cref{thm:p03-w03-p03-067-as-theory}.  Without completeness, item 2 fails:
the empty theory has only almost sure members, and in \(SU(n,c/n)\) the
sentence \(\exists x\,U(x)\) has limit \(1-e^{-c}\)
(\cref{thm:p03-w03-p03-064-yes}).

\begin{hbmethod}{def:p03-w03-p03-072-method}{Proving a zero--one law with the game}
\begin{enumerate}
\item Write down an explicit theory \(T\) in a language with equality and
  finitely many relation symbols.
\item Show that every member of \(T\) holds almost surely.  This is the
  probabilistic half, and it is where the regime of \(p(n)\) enters.  It
  gives the consistency of \(T\) (\cref{thm:p03-w03-p03-072-transfer},
  item 1) and, when \(T\) contains the sentences ``there are \(r\) distinct
  elements'', that \(T\) has no finite models.
\item For every \(t\geq1\) and every two countable models \(G,H\) of \(T\),
  exhibit a winning strategy for Duplicator in \(\mathrm{EHR}(G,H;t)\).  A
  strategy that works for models of any cardinality may be used instead.
\item Conclude that \(T\) is complete, by
  \cref{thm:p03-w03-p03-064-complete-game} in the first case and, in the
  second, because any two models of \(T\) are then elementarily equivalent
  by \cref{thm:p03-w03-p03-072-game-theorem}.  Conclude the zero--one law by
  \cref{thm:p03-w03-p03-072-transfer}, item 2.
\end{enumerate}
Side conditions.  Step 3 needs a strategy for each finite \(t\); a
strategy in the game of length \(\omega\) is more than is needed, and on
countable models it gives an isomorphism
(\cref{thm:p03-w03-p03-064-infinite-game}).  Step 3 must be carried out on
models of \(T\), not on the finite random structures, which satisfy each
axiom only with high probability.  The limit obtained is \(1\) exactly for
the consequences of \(T\); so if one theory \(T\) serves for a whole range
of functions \(p(n)\), the limit of each sentence is the same throughout
that range.
\end{hbmethod}

Spencer carries out the method three times.  For the unary model
\(SU(n,p)\) with \(np(n)\to\infty\) and \(n(1-p(n))\to\infty\), Duplicator
answers a new element satisfying \(U\) by a new element satisfying \(U\),
and likewise for \(\neg U\)
\cite[p.~3]{Spencer2007LogicRandomStructuresAuthor}; this is
\cref{thm:p03-w03-p03-067-unary}.  For \(G(n,p)\) with \(p\) constant,
Duplicator stays alive using the extension axioms
\cite[pp.~12--13]{Spencer2007LogicRandomStructuresAuthor}; this is
\cref{thm:p03-w03-p03-064-alice-complete,thm:p03-w03-p03-067-complete}.  For
\(p=n^{-\alpha}\) with \(\alpha\) irrational, staying alive is not enough
and Duplicator must look ahead; this is the next section.

\section{The closure strategy as a winning strategy}
\label{sec:p03-w03-p03-072-closure}

In \(G(n,n^{-\alpha})\) a set of vertices may have a special property, for
instance three vertices having a common neighbour, which a random set of
the same size almost surely lacks; such properties are given by rigid
extensions \cite[pp.~7--8]{Spencer2007LogicRandomStructuresAuthor}.
Duplicator copes by keeping not only the chosen vertices but their closures
under small rigid extensions isomorphic, and by looking further ahead the
more rounds remain \cite[pp.~11--12]{Spencer2007LogicRandomStructuresAuthor}.
The sequence of look-ahead bounds is the following.

\begin{hbdefinition}{def:p03-w03-p03-072-lookahead}{The look-ahead bounds}
Let \(\lambda\in(0,1]\) and let \(u\geq1\) be the number of rounds.  The
\emph{look-ahead bounds} for \((\lambda,u)\) are the strict sequence
\(s^\lambda_0,\ldots,s^\lambda_u\) of \cref{sec:p03-w03-p03-070-strategy}:
\(s^\lambda_0=0\), \(s^\lambda_1=1\) and, for \(1\leq j\leq u-1\),
\[
  s^\lambda_{j+1}=\max\Bigl(s^\lambda_j,\
  \Bigl\lfloor\frac{u-j}{\varepsilon_\lambda(s^\lambda_j)}\Bigr\rfloor+1\Bigr),
\]
with \(\varepsilon_\lambda\) as defined after
\cref{thm:p03-w03-p03-070-closure}.  With \(i\) rounds still to be played,
Duplicator works with closures under rigid extensions of at most
\(s^\lambda_i\) vertices.
\end{hbdefinition}

Spencer chooses \(t_0=0\), \(t_1=1\) and each \(t_{i+1}\geq t_i\) so large
that almost surely the \(t_i\)-closure of any \(i+1\) vertices has at most
\(t_{i+1}\) vertices outside them, which the Finite Closure Lemma makes
possible \cite[p.~11]{Spencer2007LogicRandomStructuresAuthor}.  The
sequence above meets that requirement in every model of the axioms, not
only almost surely, by the counting form of the lemma
(\cref{thm:p03-w03-p03-070-finite-closure}); it is indexed by the number of
rounds remaining.

\begin{hbtheorem}{thm:p03-w03-p03-072-closure-strategy}{proved=annex:p03-w03-p03-072-closure-strategy}{Duplicator's closure strategy}
Let \(\lambda\in(0,1]\), let \(u\geq1\), and let \(G\) and \(G'\) be models
of \(T^-_\lambda\) (\cref{def:p03-w03-p03-070-theory}) of any cardinality;
for irrational \(\alpha\in(0,1)\), \(T^-_\alpha=T_\alpha\), the theory of
the Nonexistence and Generic Extension schemas.  Let
\(t_i=s^\lambda_i\) be the look-ahead bounds, closures being taken at
\(\lambda\) in both graphs.
\begin{enumerate}
\item Duplicator has a strategy in \(\mathrm{EHR}(G,G';u)\) that keeps, for
  each \(i=u,u-1,\ldots,0\), the following invariant after round
  \(s=u-i\): there is an isomorphism
  \[
    \varphi\colon G|_{\operatorname{cl}_{t_i}(x_1,\ldots,x_s)}\to
    G'|_{\operatorname{cl}_{t_i}(x'_1,\ldots,x'_s)}
  \]
  with \(\varphi(x_j)=x'_j\) for \(j\leq s\), where \(x_j\in G\) and
  \(x'_j\in G'\) are the vertices chosen in round \(j\), whoever chose them.
\item Every play following this strategy is won by Duplicator.
\item Any two models of \(T^-_\lambda\) are elementarily equivalent, and
  \(T^-_\lambda\) is complete if it is consistent.
\end{enumerate}
\end{hbtheorem}

At the start the invariant says that the \(t_u\)-closures of the empty set
agree, which the Nonexistence schema secures by making them empty; at the
end, with \(t_0=0\), it says that \(x_j\mapsto x'_j\) is an isomorphism of
the chosen vertices, which is Duplicator's winning condition.  Each round
follows Spencer's two cases.  If Spoiler's vertex \(y\) lies inside the old
closure \(\mathrm{OLD}=\operatorname{cl}_{\mathrm{BIG}}(x_1,\ldots,x_s)\),
with \(\mathrm{BIG}=t_i\), Duplicator answers with its image under
\(\varphi\) (Inside).  Otherwise the part \(\mathrm{NEW}\) of
\(\operatorname{cl}_{\mathrm{SMALL}}(x_1,\ldots,x_s,y)\), with
\(\mathrm{SMALL}=t_{i-1}\), that lies outside \(\mathrm{OLD}\) is safe over
\(\mathrm{OLD}\), and Duplicator copies it over the image of \(\mathrm{OLD}\)
by the Generic Extension schema with genericity parameter \(\mathrm{SMALL}\),
answering with the copy of \(y\) (Outside)
\cite[p.~12]{Spencer2007LogicRandomStructuresAuthor}.  Spencer leaves out
the lemma that a nonsafe extension has a rigid subextension and the
genericity lemmas that make the new \(\mathrm{SMALL}\)-closures isomorphic
\cite[p.~12]{Spencer2007LogicRandomStructuresAuthor}; they are
\cref{thm:p03-w03-p03-070-weights}, item 3, and
\cref{thm:p03-w03-p03-070-generic-closure}, and the whole round is
\cref{thm:p03-w03-p03-070-extension}.

\begin{hbcorollary}{thm:p03-w03-p03-072-sparse-zero-one}{proved=annex:p03-w03-p03-072-sparse-zero-one}{The zero--one law for \(n^{-\alpha}\) by the game}
Let \(\alpha\in(0,1)\) be irrational.  Then \(T_\alpha\) of
\cref{def:p03-w03-p03-070-theory} is consistent and complete, every
first-order graph sentence \(B\) satisfies
\(\lim_{n\to\infty}\Pr[G(n,n^{-\alpha})\entails B]\in\{0,1\}\), and the
limit is \(1\) iff \(T_\alpha\entails B\).
\end{hbcorollary}

This is the conclusion of \cref{thm:p03-w03-p03-070-complete}, item 2, now
reached through the game, which is Spencer's route
\cite[p.~12]{Spencer2007LogicRandomStructuresAuthor}.  The proved input
\((\mathrm{AS}_\alpha)\) of \cref{thm:p03-w03-p03-070-as-holds} supplies
the probabilistic half.  For the Nonexistence
schema it is the first-moment count of
\cref{thm:p03-w03-p03-068-nonexistence}.  For Generic Extension Spencer
derives it from a large deviation bound for the number of extensions and a
counting argument, and omits the technical part
\cite[pp.~10--11]{Spencer2007LogicRandomStructuresAuthor}; the handbook
proves the counting estimate and the Generic Extension schema almost surely
as \cref{thm:p03-w03-p03-065-counting,thm:p03-w03-p03-065-generic-extension}.
These estimates establish \((\mathrm{AS}_\alpha)\); the annex combines
it with the closure strategy and the transfer theorem to obtain the
Shelah--Spencer law \cref{thm:p03-w03-p03-065-shelah-spencer}.

\begin{hbexample}{def:p03-w03-p03-072-pi-seven}{Two rounds at \(\alpha=\pi/7\)}
Take Spencer's \(\alpha=\pi/7=0.4487\ldots\)
\cite[p.~7]{Spencer2007LogicRandomStructuresAuthor} and \(u=2\).  For
\(v=1\) the least \(e\) with \(e\alpha>1\) is \(e=3\), since
\(2\alpha=0.897\ldots<1<3\alpha=1.346\ldots\); so
\(\varepsilon_\alpha(1)=3\alpha-1=0.346\ldots\), and the look-ahead bounds
are \(t_0=0\), \(t_1=1\) and
\(t_2=\max(1,\lfloor1/0.346\ldots\rfloor+1)=\max(1,2+1)=3\).  A single vertex
\(z\) is rigid over a set \(X\) exactly when it has at least three
neighbours in \(X\), because \(1-e\alpha<0\) iff \(e\geq3\).
\begin{itemize}
\item Round 1 (\(i=2\), \(\mathrm{BIG}=3\), \(\mathrm{SMALL}=1\)): by
  Nonexistence, \(\mathrm{OLD}=\operatorname{cl}_3(\varnothing)=\varnothing\)
  in both graphs, so Spoiler's \(x_1\) is Outside.  No vertex has three
  neighbours in \(\{x_1\}\), so \(\operatorname{cl}_1(x_1)=\{x_1\}\) and
  \(\mathrm{NEW}=\{x_1\}\), with \(v=1\), \(e=0\): safe over
  \(\varnothing\).  For the same reason every vertex is \(1\)-generic over
  the empty tuple, and Duplicator may answer any vertex \(x'_1\).
\item Round 2 (\(i=1\), \(\mathrm{BIG}=1\), \(\mathrm{SMALL}=0\)):
  \(\mathrm{OLD}=\operatorname{cl}_1(x_1)=\{x_1\}\).  If Spoiler repeats
  \(x_1\), the move is Inside and Duplicator repeats \(x'_1\).  Otherwise
  \(\mathrm{NEW}=\{x_2\}\) with \(v=1\), \(e\in\{0,1\}\); since
  \(1-\alpha>0\) it is safe, and Duplicator answers with a vertex
  \(x'_2\neq x'_1\) adjacent to \(x'_1\) iff \(x_2\sim x_1\), which Generic
  Extension provides.
\end{itemize}
In two rounds the strategy reduces to matching adjacency.  The look-ahead
matters in longer games.  If \(x_1,x_2,x_3\) have been chosen with rounds
still to come and they have a common neighbour \(z\), then \(\{z\}\) is
rigid over \(\{x_1,x_2,x_3\}\) and lies in their \(t\)-closure for every
\(t\geq1\).  Since Duplicator has kept those closures isomorphic,
\(x'_1,x'_2,x'_3\) have a common neighbour as well, and a Spoiler move to
\(z\) is Inside.  Matching only the adjacencies among the chosen vertices
does not secure this: three vertices of \(G(n,n^{-\pi/7})\) chosen at random
almost surely have no common neighbour
\cite[pp.~7--8]{Spencer2007LogicRandomStructuresAuthor}, so an answer
triple chosen by adjacency alone need not have one.
\end{hbexample}

\begin{hbwarning}{def:p03-w03-p03-072-limits}{What the closure strategy does not give}
The strategy proves that the axioms decide every sentence; it says nothing
about probabilities.  The zero--one law needs, in addition, that every axiom
holds almost surely, a separate probabilistic fact
(\cref{thm:p03-w03-p03-065-generic-extension,thm:p03-w03-p03-068-nonexistence}).
\Cref{thm:p03-w03-p03-070-as-holds} establishes that fact for the irrational
exponents used in \cref{thm:p03-w03-p03-072-sparse-zero-one}.  The
strategy is played on models of all the axioms, which are infinite, since
the Generic Extension sentence for \(v\) nonroots without roots or edges
asks for \(v\) distinct vertices; a finite random graph satisfies any fixed
finite set of axioms only with high probability.  For rational \(\lambda=a/b\) the theory \(T^-_\lambda\), which
counts weight-zero extensions as sparse, is still complete when consistent
by item 3 of \cref{thm:p03-w03-p03-072-closure-strategy}, but completeness
there is not a zero--one law at \(\lambda\).  Spencer remarks that the
zero--one law is not expected at rational exponents
\cite[p.~19]{Spencer2007LogicRandomStructuresAuthor};
\cref{def:p03-w03-p03-065-rational-boundary} shows an extension statement
whose behaviour changes at such a point; and what \(T^-_\lambda\) decides is
the left limit of \cref{thm:p03-w03-p03-070-left-limit}.  The
strategy also does not give a strategy for the game of length \(\omega\):
the bounds \(t_i\) depend on the number \(u\) of rounds, and a countable
model of \(T_\alpha\) need not be unique up to isomorphism
(\cref{thm:p03-w03-p03-067-sparse-models}).
\end{hbwarning}

\section{Relationships}
\label{sec:p03-w03-p03-072-relationships}

The game is the finite-relational case of \cref{def:p03-ef-game}, and the
game theorem is \cref{thm:p03-w03-p03-038-pointed-ef} without parameters.
The earlier games of \cref{sec:p03-ef-games} compare pairs of finite
structures, such as finite linear orders of different lengths, to show that
a property is not first-order; here the game compares arbitrary models of
an almost sure theory to show that the theory is complete.  The constant
and unary applications are
\cref{thm:p03-w03-p03-064-alice-complete,thm:p03-w03-p03-067-complete,thm:p03-w03-p03-067-unary};
Fagin's route to the fixed-\(p\) graph law is \cref{thm:p09-fagin1976-graph};
the uniform-distribution zero--one law of \cref{sec:p03-zero-one-law} is
the case \(p=1/2\).  The rooted-graph notions and the almost sure behaviour
of extension statements are in \cref{sec:p03-w03-p03-065-extension-statements}.
The same one-round step, run between models of \(T^-_\lambda\) and
\(T^-_\mu\) for nearby parameters, gives the Jump and Complexity Conditions
(\cref{thm:p03-w03-p03-070-transfer}); for rational exponents Spencer also
describes a different game, played on closure types level by level rather
than on vertices \cite[pp.~20--21]{Spencer2007LogicRandomStructuresAuthor},
which is \cref{def:p03-w03-p03-070-game}.

\section{Sources}
\label{sec:p03-w03-p03-072-sources}

Spencer, ``Logic and Random Structures'', the version on his web page
\cite{Spencer2007LogicRandomStructuresAuthor}: almost sure and almost never,
p.~2; the reduction of a zero--one law to a complete almost sure theory and
the unary strategy, p.~3; rooted graphs and the examples at
\(\alpha=\pi/7\), pp.~6--8; the almost sure theory \(T_\alpha\), pp.~9--11;
the closure strategy, pp.~11--12; the constant-\(p\) strategy,
pp.~12--13; the sentences ``there are \(r\) distinct elements'' in an almost
sure theory, p.~13; the level game for rational exponents, pp.~20--21; the
Ehrenfeucht game, its theorem, the isolated-vertex example, the two
corollaries and the extension to finitely many relation symbols,
pp.~26--27.  Spencer attributes the game to Ehrenfeucht, with an earlier
form in Fra\"iss\'e's work, and the sparse zero--one law to Shelah and
himself \cite[p.~27]{Spencer2007LogicRandomStructuresAuthor}.
